% Source: https://artofproblemsolving.com/community/c3h2038778_geometry_problem_need_help
% Stable identity: aops-topic-2038778-q01
Let $r$ be the radius. The circle in the corner at $A$ is tangent to both sides at $A$, so its centre $O_1$ is on the bisector of the $60^\circ$ angle at $A$, at height $r$ above $AB$. In the right triangle formed by $A$, $O_1$ and the foot of the perpendicular from $O_1$ to $AB$, the angle at $A$ is $30^\circ$, so the foot is at distance $r\sqrt3$ from $A$. The same holds at $B$.

\begin{asy}
usepackage("amsmath");
size(6cm);
defaultpen(fontsize(10pt));
real r=sqrt(3); real s=6+2*sqrt(3);
pair A=(0,0), B=(s,0), C=(s/2,s*sqrt(3)/2);
pair O1=(r*sqrt(3),r), O2=(s-r*sqrt(3),r), O3=(s/2,s*sqrt(3)/2-2*r);
draw(A--B--C--cycle,linewidth(0.9));
draw(circle(O1,r),blue); draw(circle(O2,r),blue); draw(circle(O3,r),blue);
draw(O1--O2,red+dashed); draw(O1--(O1.x,0),red); draw(A--O1,gray+dashed);
dot(O1);dot(O2);dot(O3);
label("$A$",A,SW); label("$B$",B,SE); label("$C$",C,N);
label("$O_1$",O1,NW); label("$O_2$",O2,NE); label("$O_3$",O3,N);
label("$r$",(O1.x,r/2),E,red); label("$2r$",O1+0.3*(O2-O1),N,red);
label("$r\sqrt3$",(r*sqrt(3)/2,0),S,fontsize(8pt)); label("$r\sqrt3$",(s-r*sqrt(3)/2,0),S,fontsize(8pt));
\end{asy}

The centres $O_1$ and $O_2$ are both at height $r$ above $AB$, so $O_1O_2=AB-2r\sqrt3$, and tangency of the two circles means $O_1O_2=2r$. Hence
\[6+2\sqrt3=2r+2r\sqrt3=2r\left(1+\sqrt3\right)\quad\Longrightarrow\quad r=\dfrac{3+\sqrt3}{1+\sqrt3}=\dfrac{\sqrt3\left(\sqrt3+1\right)}{1+\sqrt3}=\sqrt3.\]
By symmetry the same $r$ works for the other two pairs. The total area is
\[3\cdot\pi\left(\sqrt3\right)^2=\boxed{9\pi}.\]
